\chapter{SQL -- Einführung in Datenbankabfragen}
\renewcommand{\chaptertitle}{SQL -- Einführung in Datenbankabfragen}

\lehead[]{\normalfont\sffamily\hspace*{-2.00cm}\textcolor{white}{\colorbox{lightblue}{\makebox[1.60cm][r]{\thechapter}}}\hspace{0.17cm}\textcolor{lightblue}{\chaptertitle}}
\rohead[]{\textcolor{lightblue}{\chaptertitle}\normalfont\sffamily\hspace*{0.17cm}\textcolor{white}{\colorbox{lightblue}{\makebox[1.60cm][l]{\thechapter}}}\hspace{-2.00cm}}
%\chead[]{}
\rehead[]{\textcolor{lightblue}{AvHG, Inf, My}}
\lohead[]{\textcolor{lightblue}{AvHG, Inf, My}}

\lstset{style=mySQL}

\section{Die Beispieldatenbank \texttt{haustier}}

In diesem Kapitel konzentrieren wir uns zunächst vollständig darauf, Daten
gezielt aus einer bestehenden Datenbank auszulesen (\lstinline|SELECT|-Abfragen).
Wie man eigene Datenbanken und Tabellen von Grund auf selbst entwirft und
verändert, lernst du im darauffolgenden Kapitel.

Für alle Übungen arbeiten wir mit der fertigen Datenbank \myFile{haustier}. Sie
verwaltet die Tiere einer Kleintierpraxis sowie deren Besitzerinnen und
Besitzer. Da ein Besitzer mehrere Tiere haben kann und ein Tier (z.\,B.\ innerhalb
einer Familie) auch mehreren Besitzern gleichzeitig gehören kann, besteht die
Datenbank aus drei Tabellen:

\begin{compactitem}
  \item \textbf{\texttt{tier}}: Speichert alle Tiere mit \lstinline|tier_id|
  (Primärschlüssel), \lstinline|name|, \lstinline|tierart|, \lstinline|lebendig|
  (\lstinline|'ja'|/\lstinline|'nein'|), \lstinline|geschlecht|,
  \lstinline|geburtstag| und \lstinline|todestag|.
  \item \textbf{\texttt{besitzer}}: Speichert alle Personen bzw.\ Einrichtungen
  mit \lstinline|besitzer_id| (Primärschlüssel), \lstinline|anrede|,
  \lstinline|vorname|, \lstinline|nachname|, \lstinline|straße|, \lstinline|plz|,
  \lstinline|ort| und \lstinline|telefonnr|.
  \item \textbf{\texttt{beziehung}}: Verknüpft über die beiden Fremdschlüssel
  \lstinline|beziehung_besitzer_id| und \lstinline|beziehung_tier_id| die
  Besitzer mit ihren Tieren ($n:m$-Beziehung).
\end{compactitem}

\noindent \textbf{Vorbereitung:} Lade die Datei \myFile{haustier.sql} aus dem
Kurs-Repository, führe das Skript in DBeaver (bzw.\ MySQL) aus und aktiviere die
Datenbank für deine Abfragen mit:

\begin{lstlisting}
USE haustier;
\end{lstlisting}

\noindent \emph{Tipp zur Schreibweise:} SQL-Schlüsselwörter sind nicht
groß-/kleinschreibungsgebunden. Gewöhne dir dennoch an, SQL-Schlüsselwörter in
GROSSBUCHSTABEN und Tabellen- sowie Spaltennamen in Kleinbuchstaben zu
schreiben. Zeichenketten sowie Datumsangaben (\lstinline|'YYYY-MM-DD'|) stehen
immer in einfachen Anführungszeichen.


\section{Grundlagen der \texttt{SELECT}-Abfrage}

Das Auslesen von Informationen aus Tabellen geschieht mit der
\lstinline|SELECT|-Anweisung:

\begin{lstlisting}
SELECT `\slshape{spalte1, spalte2, \ldots}`
FROM `\slshape{tabelle1, tabelle2, \ldots}`
[WHERE `\slshape{Bedingung}`]
[GROUP BY `\slshape{gruppe}`]
[HAVING `\slshape{Gruppen\_Bedingung}`]
[ORDER BY `\slshape{sortierspalten}`];
\end{lstlisting}

Hinter dem Schlüsselwort \lstinline|SELECT| gibt man einen oder mehrere
Spaltennamen an, die ausgelesen werden sollen. Möchte man alle Spalten einer
Tabelle anzeigen, schreibt man \lstinline|*|. Müssen Spalten aus verschiedenen
Tabellen unterschieden werden, nutzt man die Punktnotation
\texttt{\slshape{tabellenname.spaltenname}}.

\subsection{Filtern mit \texttt{WHERE} und Vergleichsoperatoren}

Hinter \lstinline|WHERE| legt man fest, welche Bedingung die ausgewählten Zeilen
erfüllen müssen. Dazu stehen in SQL folgende Vergleichs- und Logikoperatoren zur
Verfügung:

\begin{table}[h]
\centering
\begin{tabular}{|c|l|}
\hline
\textbf{Operator} & \textbf{Bedeutung}\\ \hline
\texttt{=} & gleich \\ \hline
\texttt{!=} oder \texttt{<>} & ungleich \\ \hline
\texttt{<} / \texttt{>} & kleiner / größer \\ \hline
\texttt{<=} / \texttt{>=} & kleiner oder gleich / größer oder gleich \\ \hline
\texttt{x IS NULL} & testet, ob \texttt{x} leer (\texttt{NULL}) ist (nicht \texttt{x = NULL}!) \\ \hline
\texttt{x IS NOT NULL} & testet, ob \texttt{x} einen Wert besitzt \\ \hline
\texttt{x IN (w1, w2, ...)} & testet, ob \texttt{x} in der angegebenen Werteliste liegt \\ \hline
\texttt{AND} / \texttt{OR} / \texttt{NOT} & logisches Und / Oder / Nicht \\ \hline
\end{tabular}
\caption{Wichtige SQL-Operatoren für Bedingungen}
\label{tab:sql-vergleichsoperatoren}
\end{table}

\subsection{Duplikate entfernen mit \texttt{DISTINCT}}

Mit dem Schlüsselwort \lstinline|DISTINCT| direkt hinter \lstinline|SELECT|
werden mehrfach vorkommende Zeilen aus dem Abfrageergebnis entfernt:

\begin{lstlisting}
SELECT DISTINCT tierart FROM tier;
\end{lstlisting}

\subsection{Suchergebnisse sortieren mit \texttt{ORDER BY}}

Mit \lstinline|ORDER BY| werden die Suchergebnisse nach einer oder mehreren
Spalten sortiert -- wahlweise aufsteigend (\lstinline|ASC|, Standardeinstellung)
oder absteigend (\lstinline|DESC|):

\begin{lstlisting}
SELECT name, tierart
FROM tier
ORDER BY tierart ASC, name DESC;
\end{lstlisting}

\subsection{Aggregatfunktionen, \texttt{GROUP BY} und \texttt{HAVING}}

Mit Aggregatfunktionen lassen sich Werte über mehrere Zeilen hinweg
zusammenfassen:

\begin{compactitem}
  \item \lstinline|COUNT(*)| zählt alle Zeilen einer Tabelle bzw.\ Gruppe;
  \lstinline|COUNT(spalte)| zählt nur Zeilen, in denen die Spalte nicht
  \lstinline|NULL| ist.
  \item \lstinline|SUM(spalte)|, \lstinline|AVG(spalte)|, \lstinline|MIN(spalte)|
  und \lstinline|MAX(spalte)| berechnen Summe, Durchschnitt, Minimum bzw.\
  Maximum einer Spalte.
\end{compactitem}

\noindent \textbf{Achtung:} Vor der öffnenden Klammer einer Aggregatfunktion
darf kein Leerzeichen stehen! Mit \lstinline|AS "Spaltenname"| kann man einer
berechneten Spalte zudem eine aussagekräftige Überschrift geben.

Mit \lstinline|GROUP BY| werden Zeilen mit gleichen Werten zu Gruppen
zusammengefasst. Möchte man anschließend ganze Gruppen nach einer Bedingung
filtern, verwendet man \lstinline|HAVING| (während \lstinline|WHERE| einzelne
Zeilen \emph{vor} der Gruppenbildung filtert):

\begin{lstlisting}
SELECT tierart, COUNT(*) AS "Anzahl"
FROM tier
GROUP BY tierart
HAVING COUNT(*) >= 2;
\end{lstlisting}
